\section*{अध्याय \ref{ch:b3:field-lab-methods-3} --- क्षेत्र और प्रयोगशाला विधियाँ III}
% Values from the Book 5 scratch calculator calc/ch31.py.

\begin{solution}{exo:b3:field-lab-methods-3:1}
(क)~$\qty{0.01}{\milli\gal} = \qty{e-7}{m.s^{-2}}$, और
$\num{e-7}/\num{9.81} = \num{1.0e-8}$: दस करोड़ में एक भाग।

(ख)~$0.01/0.3086 = \qty{0.032}{m}$, लगभग \qty{3.2}{cm}।

(ग)~दोनों प्रभाव मिलकर गुरुत्व को
\qty{0.1967}{\milli\gal.m^{-1}} से बदलते हैं, जो अंतर $0.3086 - 0.1119$ है, इसलिए
ऊँचाई में $0.01/0.1967 = \qty{5.1}{cm}$ का बदलाव होना चाहिए।
\end{solution}

\begin{solution}{exo:b3:field-lab-methods-3:2}
$\Delta m = \tfrac43\pi(\qty{50}{m})^3\times\qty{500}{kg.m^{-3}} =
\qty{2.62e8}{kg}$। शिखर-मान $G\Delta m/z^2 =
\num{6.674e-11}\times\num{2.62e8}/\num{e4} = \qty{1.75e-6}{m.s^{-2}}$,
अर्थात \qty{0.175}{\milli\gal} है। अर्धचौड़ाई $x_{1/2} = z/1.305 =
\qty{77}{m}$ है।
\end{solution}

\begin{solution}{exo:b3:field-lab-methods-3:3}
(क)~$AM = MN = NB = a$ मानने पर, $K = 2\pi/(1/a - 1/2a - 1/2a + 1/a) =
2\pi a = \qty{62.83}{m}$, और $\rho_a = 62.83\times0.239/0.1 =
\qty{150}{\ohm\metre}$। (ख)~$K = \pi(15^2 - 5^2)/(2\times5) = \pi\times
200/10 = \qty{62.83}{m}$। तब इलेक्ट्रोड $0$, $10$, $20$
और \qty{30}{m} पर होते हैं, AMNB क्रम में; यह
\qty{10}{m} अंतराल वाला वेनर विन्यास है।
\end{solution}

\begin{solution}{exo:b3:field-lab-methods-3:4}
(क)~$V = 0.2998/\sqrt5 = \qty{0.134}{m.ns^{-1}}$। (ख)~$t =
2\times10/0.134 = \qty{149}{ns}$। (ग)~$\delta = 0.0053\sqrt5/\num{e-3} =
\qty{11.9}{m}$, और लगभग \qty{36}{m} की उपयोगी गहराई। (घ)~मिट्टी में
$\delta = 0.0053\sqrt{20}/0.01 = \qty{2.4}{m}$, लगभग
\qty{7}{m} की उपयोगी गहराई, और केवल \qty{0.067}{m.ns^{-1}} का वेग।
सन्निकटन वहाँ भी लागू रहता है: \qty{100}{MHz} पर,
$\omega\varepsilon = 2\pi\times\num{e8}\times20\times\num{8.85e-12} =
\qty{0.11}{S.m^{-1}}$, जो $\sigma$ का दस गुना है।
\end{solution}

\begin{solution}{exo:b3:field-lab-methods-3:5}
(क)~बेलन का $\lambda = \pi R^2\Delta\rho$ प्रति इकाई लंबाई है, और
उसके बाहर उसका आकर्षण अपने अक्ष पर स्थित रेखीय द्रव्यमान जैसा होता है (गॉस
का प्रमेय, सहअक्षीय बेलन पर): $\Delta g = 2G\pi R^2\Delta\rho\,z/(x^2
+ z^2)$। शिखर का आधा तब प्राप्त होता है जब $x^2 + z^2 = 2z^2$, इसलिए $z =
x_{1/2}$। (ख)~गोला $1.305\times100 = \qty{130}{m}$ पर है, जबकि
बेलन \qty{100}{m} पर है: गोले के रूप में पढ़ी गई वही परिच्छेदिका पिंड को
\qty{30}{\percent} अधिक गहराई पर रखती है। गोले की विसंगति अधिक तेज़ी से
दूरी के साथ क्षीण होती है, इसलिए गोले को उतना चौड़ा होने के लिए अधिक गहरा होना पड़ता है।
\end{solution}

\begin{solution}{exo:b3:field-lab-methods-3:6}
(क)~प्रत्येक बिंदु $\qty{100}{m}\times\qty{100}{m} =
\qty{e4}{m^2}$ का प्रतिनिधित्व करता है, इसलिए $\iint\Delta g\,\mathrm{d}A \approx
\num{-438e-5}\times\num{e4} = \qty{-43.8}{m^3.s^{-2}}$, और $\Delta m
= -43.8/(2\pi\times\num{6.674e-11}) = \qty{-1.04e11}{kg}$। (ख)~$1.04/1.74
= \qty{60}{\percent}$। शेष भाग विसंगति के उन पार्श्वों में है
जो मानचित्र के बाहर हैं: ग्रेनाइट एक किलोमीटर से अधिक गहराई तक पहुँचता है, और
उसके गहरे भाग की विसंगति दूरी के साथ धीरे-धीरे क्षीण होती है। मानचित्र को
पर्याप्त दूर तक बढ़ाने पर, या किसी मॉडल से, उसे पुनः प्राप्त किया जा सकता है।
\end{solution}

\begin{solution}{exo:b3:field-lab-methods-3:7}
(क)~पतले डाइक के लिए विसंगति $x = \pm z$ पर अपने शिखर के आधे तक गिर जाती है:
आधे शिखर पर इसकी पूर्ण चौड़ाई $2z = \qty{60}{m}$ है, और
इसके आर-पार तीन पाठ्यांक लेने के लिए अधिकतम \qty{20}{m} का अंतराल चाहिए। (ख)~$2\times
220/3 = \qty{147}{m}$। (ग)~आवश्यक स्टेशन-अंतराल सबसे उथले रुचि के पिंड की
गहराई के अनुपात में होता है: गहरे
डाइक के लिए, \qty{150}{m} पर, बनाया गया सर्वेक्षण उथले की विसंगति को उपनामित कर देगा, जो
केवल \qty{60}{m} चौड़ी है; उथले पिंड के लिए बनाया गया सर्वेक्षण गहरे को भी
देख पाता है, पर उसके लिए आवश्यक पाठ्यांकों से कहीं अधिक पाठ्यांक लेता है।
\end{solution}

\begin{solution}{exo:b3:field-lab-methods-3:8}
(क)~$F = 80/8 = 10$ और $\phi = F^{-1/m} = 10^{-1/2} = \num{0.316}$।
(ख)~$10^{-1/1.8} = \num{0.278}$ और $10^{-1/2.2} = \num{0.351}$।
$m$ में \num{0.2} का परिवर्तन छिद्रता को लगभग चार प्रतिशत-बिंदु खिसका देता है;
प्रतिरोधकताएँ, जो कुछ प्रतिशत तक मापी गई हैं, उसे बहुत कम खिसकाती हैं।
सीमेंटन घातांक अनिश्चितता को नियंत्रित करता है, और
इसे उसी चट्टान के नमूनों पर अंशांकित किया जाना चाहिए।
\end{solution}

\begin{solution}{exo:b3:field-lab-methods-3:9}
(क)~$s = 2h_1$ पर, $(2nh_1/s)^2 = n^2$ होने से, पद हैं $k^n(1 +
n^2)^{-3/2}$: $0.3536k$, $0.0894k^2$, $0.0316k^3$। $k = 0.5$ के लिए:
$\num{0.1768}$, $\num{0.0224}$, $\num{0.0040}$, और $\rho_a/\rho_1 \approx
1 + 2\times0.2031 = \num{1.406}$ (पूर्ण योग \num{1.409} देता है)। $k
= -0.5$ के लिए पद चिह्न में एकांतरित होते हैं, $-0.1768 + 0.0224 - 0.0040$, और
$\rho_a/\rho_1 \approx \num{0.683}$ (पूर्ण योग \num{0.685})। (ख)~$s \ll
h_1$ के लिए प्रत्येक पद लुप्त हो जाता है और $\rho_a \to \rho_1$; $s \gg h_1$ के लिए प्रत्येक
कोष्ठक~1 की ओर प्रवृत्त होता है और $1 + 2\sum k^n = 1 + 2k/(1 - k) = (1 + k)/(1 -
k) = \rho_2/\rho_1$। (ग)~$\rho_2/\rho_1 = (1 + k)/(1 - k) = 3$ और $1/3$।
\end{solution}

\begin{solution}{exo:b3:field-lab-methods-3:10}
(क)~अंतराल $a$ के साथ AMNB क्रम में: $AM = a$, $BM = 2a$, $AN = 2a$,
$BN = a$, इसलिए कोष्ठक $2/a - 1/a = 1/a$ है और $K = 2\pi a$।
(ख)~A और B को $\mp s$ पर और M तथा N को $\mp l$ पर रखने पर: $AM = BN = s - l$,
$BM = AN = s + l$, कोष्ठक है $2/(s - l) - 2/(s + l) = 4l/(s^2 -
l^2)$, और $K = \pi(s^2 - l^2)/(2l)$। जब $l \ll s$, $\Delta V_{MN}
\approx 2l\,E$ है, जहाँ $E = 2(-\mathrm{d}V/\mathrm{d}r)$ केंद्र पर,
दोनों धारा-इलेक्ट्रोडों का क्षेत्र; तब $\rho_a = K\Delta V/I \to
\pi s^2/(2l)\times 4l(-\mathrm{d}V/\mathrm{d}r)/I = 2\pi s^2(-
\mathrm{d}V/\mathrm{d}r)/I$। (ग)~सन्निकट गुणक और सटीक गुणक का अनुपात
$s^2/(s^2 - l^2) = 1/(1 - 0.01) = 1.010$ है, यानी एक प्रतिशत।
\end{solution}

\begin{solution}{exo:b3:field-lab-methods-3:11}
(क)~पाठ्यांक: \qty{0.010}{\milli\gal}। ऊँचाई:
$0.1967\times0.05 = \qty{0.010}{\milli\gal}$। घनत्व: बूगेर
सुधार $0.04191\,h\,d$ है, जहाँ $d$ \unit{g.cm^{-3}} में है, इसलिए
$0.04191\times50\times0.1 = \qty{0.210}{\milli\gal}$। भू-भाग:
\qty{0.200}{\milli\gal}। कुल: $\sqrt{0.010^2 + 0.010^2 + 0.210^2 +
0.200^2} = \qty{0.29}{\milli\gal}$। (ख)~घनत्व और भू-भाग
प्रभावी हैं, प्रत्येक पाठ्यांक-त्रुटि का लगभग बीस गुना।
\qty{0.1}{\milli\gal} का संसूचन करने के लिए, स्लैब का घनत्व नमूनों पर मापा जाए
या आँकड़ों पर समंजित किया जाए (वह मान जो स्थलाकृति को सबसे कम
प्रतिबिंबित करता है), और भू-भाग को एक सूक्ष्म उन्नयन मॉडल से संशोधित किया जाना चाहिए।
(ग)~$\sqrt{0.010^2 + 0.010^2 + 0.210^2} = \qty{0.21}{\milli\gal}$।
\end{solution}

\begin{solution}{exo:b3:field-lab-methods-3:12}
(क)~केंद्रीय पाठ्यांक से मेल खाने वाला द्रव्यमान $\Delta m = \Delta g(0)z^2/G
= \num{-4.16e-5}z^2/\num{6.674e-11}$ है: \qty{300}{m} पर \qty{-5.61e10}{kg},
\qty{400}{m} पर \qty{-9.97e10}{kg}, \qty{500}{m} पर \qty{-1.56e11}{kg}।
$x = \pm300$ और $\pm600$ पर अनुमानित पाठ्यांक
\num{-1.47} और \num{-0.37}, \num{-2.13} और \num{-0.71}, \num{-2.62}
और \qty{-1.09}{\milli\gal} हैं। \omterm{def:b3:field-lab-methods-3:inversion-practice}{बेमेल} $\chi^2 = 1238$, $3.7$ और
$856$ हैं। (ख)~केवल \qty{400}{m} वाला मॉडल, पाँच आँकड़ों के लिए $\chi^2 = 3.7$ के साथ, त्रुटियों के भीतर बैठता है।
(ग)~हाँ: कोई भी गोलाकार सममित पिंड, जिसका द्रव्यमान और केंद्र समान हो, ठीक वही विसंगति देता है,
और अन्य आकार भी त्रुटियों के भीतर आ सकते हैं। नमूनों पर मापा गया घनत्व,
कोई बोरछिद्र या कोई अन्य विधि
यह तय कर सकती है।
\end{solution}

\begin{solution}{pb:b3:field-lab-methods-3:1}
\textbf{1.} $(2846.708 - 2846.512)/\qty{2.0}{h} = \qty{0.098}{\milli\gal.h^{-1}}$.

\textbf{2.} $2851.075 - 0.098\times0.5 - 2846.512 = \qty{4.514}{\milli\gal}$;
$2850.497 - 0.088 - 2846.512 = \qty{3.897}{\milli\gal}$;
$2849.392 - 0.127 - 2846.512 = \qty{2.753}{\milli\gal}$.

\textbf{3.} बेस से ऊपर ऊँचाइयाँ $-27.1$, $-30.6$ और
\qty{-18.1}{m} हैं; $0.1967\,\Delta h$ जोड़ने पर बूगेर विसंगतियाँ
$\num{-0.817}$, $\num{-2.122}$ और \qty{-0.808}{\milli\gal} प्राप्त होती हैं।
स्टेशन आधार से नीचे हैं और पृथ्वी के केंद्र के अधिक निकट हैं,
इसलिए उनके कच्चे पाठ्यांक अधिक होते हैं; संशोधन इसे हटा देता है
और ग्रेनाइट के ऊपर एक निम्न क्षेत्र छोड़ देता है।

\textbf{4.} अतिरिक्त अंतर $0.098 - 0.040 = \qty{0.058}{\milli\gal.h^{-1}}$,
\qty{0.12}{\milli\gal} लूप पर, ठोस पृथ्वी का ज्वार है। दो
घंटों में यह लगभग रैखिक रूप से बदलता है, इसलिए दोनों आधार-पाठ्यांकों से होकर खींची सीधी रेखा
उसे विचलन के साथ हटा देती है।

\textbf{5.} त्रुटियाँ $\num{0.008}$, $\num{0.009}$ और
\qty{0.017}{\milli\gal} हैं। हर ऊँचाई को \qty{0.1}{m} तक पूर्णांकित करने से
ऊँचाई-अंतर \qty{0.1}{m} तक गलत हो सकता है, अर्थात
\qty{0.020}{\milli\gal}: पूर्णांकन त्रुटियों की व्याख्या करता है, और
सीधी रेखा के बाद बचा ज्वार \qty{0.001}{\milli\gal} से कम है।

\textbf{6.} $0.01/0.1967 = \qty{5.1}{cm}$। मैनुअलों का \qty{3.3}{cm}
$0.01/0.3086$ है: मुक्त-वायु विसंगति बूगेर सुधार के लिए चुने गए घनत्व पर
निर्भर नहीं करती, और ऊँचाई को भी इसी के लिए उपयुक्त होना चाहिए।

\textbf{7.} $x_{1/2} = (3.952 - 3.048)/2 = \qty{0.452}{km}$ और $z = 1.305
\times452 = \qty{590}{m}$ भूमि के नीचे।

\textbf{8.} $\Delta m = \Delta g\,z^2/G = \num{-2.293e-5}\times590^2/
\num{6.674e-11} = \qty{-1.20e11}{kg}$। $\Delta\rho =
\qty{-130}{kg.m^{-3}}$ लेने पर इसका आयतन \qty{9.2e8}{m^3} है और त्रिज्या
\qty{603}{m} है, जो उसके केंद्र की गहराई से बड़ी है: गोला भूमि से
बाहर निकल जाएगा। ग्रेनाइट के घनत्व-अंतर वाला कोई भी गोला
उपयुक्त नहीं बैठता; पिंड को सतह के निकट तक पहुँचना होगा।

\textbf{9.} ग्रेनाइट सतह से काफ़ी गहराई तक फैला है। इसका उथला
भाग, स्टेशनों के निकट, शिखर को सँकरा बनाता है, और नियम, जो
केवल चौड़ाई जानता है, एक ही गोले को मध्यम गहराई पर रखता है;
गहरा भाग चौड़ी ढलानें जोड़ता है, जिन्हें गोला नहीं दर्शा पाता (प्रोफ़ाइल का चित्र),
और उनमें समाहित द्रव्यमान छूट जाता है।

\textbf{10.} $2\pi G\Delta\rho = 2\pi\times\num{6.674e-11}\times(-130) =
\qty{-5.45e-8}{s^{-2}}$। जहाँ $z_1 = 0$, $z_2 = \qty{1645}{m}$, $a =
\qty{461}{m}$ लेने पर कोष्ठक का मान है $1645 + 461 - 1708.4 = \qty{397.6}{m}$,
और $\Delta g = \qty{-2.17}{\milli\gal}$ है। असीमित गहराई के लिए कोष्ठक
$a$ की ओर प्रवृत्त होता है, और $\Delta g = \qty{-2.51}{\milli\gal}$ है।

\textbf{11.} कोष्ठक $2.29/\num{5.45e-3} = \qty{420}{m}$ होना चाहिए, जिसे लगभग
\qty{2.6}{km} पर स्थित एक तल प्रदान करता है। \qty{1.6}{km} और असीमित गहराई के बीच का कोई भी तल
शिखर को केवल कुछ दशमांश मिलीगल से बदलता है,
और थोड़ा चौड़ा या घना पिंड, उथले तल के साथ, वही शिखर देता है:
तल का निर्धारण नहीं हो पाता।

\textbf{12.} $z_1 = \qty{6}{m}$ लेने पर: $1639 + 461.0 - 1708.4 = 391.6$, और
$\Delta g = \qty{-2.14}{\milli\gal}$ है, जो शीर्ष भूमि पर होने पर प्राप्त मान से
\qty{0.03}{\milli\gal} कम है। वह सर्वेक्षण की भू-भाग और घनत्व त्रुटियों से
नीचे है: गुरुत्व यह नहीं बता सकता कि ग्रेनाइट जलोढ़ के नीचे सतह तक पहुँचता है या नहीं।
ध्वनन यह बता सकता है।

\textbf{13.} $K = \pi(20^2 - 1^2)/(2\times1) = \qty{626.7}{m}$ और $\Delta V =
\rho_a I/K = 821\times0.050/626.7 = \qty{65.5}{mV}$.

\textbf{14.} $V = 2\times1/\sqrt{17.1^2 - 10.7^2} = 2/13.34 =
\qty{0.150}{m.ns^{-1}}$, $\varepsilon_r = (0.2998/0.150)^2 = 4.0$, शुष्क रेत;
पाइप की गहराई $0.150\times10.7/2 = \qty{0.80}{m}$ है।

\textbf{15.} $0.150\times20.0/2 = \qty{1.50}{m}$.

\textbf{16.} $20/821 = \num{0.0244}$, $30/1309 = \num{0.0229}$, $40/1684 =
\num{0.0238}$: $S \approx \qty{0.0237}{S}$.

\textbf{17.} शुष्क परत का योगदान $1.5/600 = \qty{0.0025}{S}$ है,
और संतृप्त परत का $\qty{0.0212}{S}$; इसलिए उसकी मोटाई $0.0212\times222 =
\qty{4.7}{m}$ है, और ग्रेनाइट $1.5 + 4.7 = \qty{6.2}{m}$
तल के नीचे है। अनुमान ऊँचा है: ग्रेनाइट असीमित रूप से
प्रतिरोधी नहीं है, इसलिए इन अंतरालों पर $\rho_a$ अब भी रेखा
$s/S$ के नीचे रहता है, और $s/\rho_a$ $S$ को अधिक आँकता है। (क्षेत्र के मॉडल में
\qty{6}{m} जलोढ़ है।)

\textbf{18.} $\phi = (20/222)^{1/2} = \num{0.30}$। $m = 1.3$ लेने पर, $\phi =
0.0901^{1/1.3} = \num{0.16}$; $m = 2.8$ लेने पर, $\phi = \num{0.42}$।

\textbf{19.} रेखा से आधार घटाने पर पश्चिम से पूर्व की ओर ये मान मिलते हैं:
3.4, 4.1, 4.8, 5.9, 7.4, 9.2, 11.5, 14.2, 17.2, 19.8, 21.0, 20.1, 17.6,
14.5, 11.5, 9.1, 7.2, 5.7, 4.6, 3.8 और \qty{3.1}{nT}।

\textbf{20.} सुधार किए बिना शिखर \qty{2.20}{km} पर \qty{28.0}{nT} है, और
इसका आधा, \qty{14.0}{nT}, पूर्व की ओर कभी नहीं पहुँचता: रेखा
\qty{15.2}{nT} पर समाप्त होती है, क्योंकि चलने के दौरान क्षेत्र
\qty{12}{nT} बढ़ गया था।

\textbf{21.} शिखर \qty{2.20}{km} पर \qty{21.0}{nT} है; इसका आधा
\qty{10.5}{nT} है। पश्चिम: \num{1.95} पर \num{9.2} और \qty{2.00}{km} पर \num{11.5} के बीच,
$1.95 + 0.05\times1.3/2.3 = \qty{1.978}{km}$ पर। पूर्व:
\num{2.40} पर \num{11.5} और \qty{2.45}{km} पर \num{9.1} के बीच,
$2.40 + 0.05\times1.0/2.4 = \qty{2.421}{km}$ पर। अर्ध-चौड़ाई है $(2.421
- 1.978)/2 = \qty{0.221}{km}$।

\textbf{22.} \ang{65} पर झुके क्षेत्र में सममित विसंगति का अर्थ है कि केवल
ऊर्ध्वाधर चुंबकन ही कार्य करता है: पिंड चुंबकीय याम्योत्तर के
अनुदिश, यानी उत्तर--दक्षिण, स्ट्राइक करता है।
\cref{prop:b3:field-lab-methods-3:dyke} का नियम एक पतली, ऊर्ध्वाधर,
द्विविमीय डाइक मानता है, जो बहुत गहराई तक पहुँचती है, प्रेरण द्वारा चुंबकित है,
और जिसका शीर्ष चुंबकीय आवेश की एक रेखा की भाँति व्यवहार करता है।

\textbf{23.} $\kappa t = 2\pi z\,\Delta T_\mathrm{max}/(F\sin^2 I) = 2\pi\times
221\times21.0/(48000\times0.8214) = \qty{0.74}{m}$। $\kappa =
0.03$ लेने पर, $t = \qty{25}{m}$। केवल गुणनफल निश्चित है: दोगुनी
मोटाई और आधी चुंबकीय प्रवृत्ति वाली डाइक वही विसंगति देती है।

\textbf{24.} विसंगति अपने शिखर के आधे पर \qty{443}{m} चौड़ी है, और स्टेशन
\qty{50}{m} दूर हैं: उसके पार नौ पाठ्यांक, जो नियम के
न्यूनतम मान से तीन गुना हैं।

\textbf{25.} डाइक का शीर्ष
\textbf{\qty[text-series-to-math=true]{220}{m}} भूमि के नीचे है। भूवैज्ञानिक
मॉडल इस बिंदु पर असंगति को \qty{222}{m} नीचे रखता है:
डाइक आवरण के आधार तक उठती है और वहीं कट जाती है, वह
उसे ढकने वाले ट्रायसिक बलुआ पत्थर से पुरानी है।
\end{solution}
